% Cuadro1
\begin{table}[htbp]
  \centering
  \caption{Estadísticas Descriptivas I}
    \begin{tabular}{rlrrrrr}
          &       & \multicolumn{4}{c}{Total}     &  \\
          &       & \multicolumn{1}{c}{media} & \multicolumn{1}{c}{desvío} & \multicolumn{1}{c}{min} & \multicolumn{1}{c}{max} & \multicolumn{1}{c}{\# true} \\
    \midrule
    \multicolumn{1}{l}{sat} & Satisfación con la vida & \multicolumn{1}{c}{1,90} & \multicolumn{1}{c}{0,87} & \multicolumn{1}{c}{1} & \multicolumn{1}{c}{4} &  \\
    \multicolumn{1}{l}{fel} & Felicidad subjetiva & \multicolumn{1}{c}{9,88} & \multicolumn{1}{c}{2,64} & \multicolumn{1}{c}{1} & \multicolumn{1}{c}{18} &  \\
    \multicolumn{1}{l}{SNU} & Uso de redes sociales & \multicolumn{1}{c}{0,79} & \multicolumn{1}{c}{0,40} & \multicolumn{1}{c}{0} & \multicolumn{1}{c}{1} & \multicolumn{1}{c}{15680} \\
    \multicolumn{1}{l}{eco\_país} & Economía del país & \multicolumn{1}{c}{3,53} & \multicolumn{1}{c}{0,93} & \multicolumn{1}{c}{1} & \multicolumn{1}{c}{5} &  \\
    \multicolumn{1}{l}{progreso} & Progreso del país & \multicolumn{1}{c}{2,10} & \multicolumn{1}{c}{0,90} & \multicolumn{1}{c}{1} & \multicolumn{1}{c}{3} &  \\
    \multicolumn{1}{l}{eco\_fliar} & Economía familiar & \multicolumn{1}{c}{2,60} & \multicolumn{1}{c}{1,10} & \multicolumn{1}{c}{1} & \multicolumn{1}{c}{5} &  \\
    \multicolumn{1}{l}{conf} & Confianza en los demás & \multicolumn{1}{c}{1,86} & \multicolumn{1}{c}{0,33} & \multicolumn{1}{c}{1} & \multicolumn{1}{c}{2} & \multicolumn{1}{c}{2486} \\
    
    \multicolumn{1}{l}{estudio} & Nivel de estudios & \multicolumn{1}{c}{1,04} & \multicolumn{1}{c}{4,32} & \multicolumn{1}{c}{1} & \multicolumn{1}{c}{17} &  \\
    \multicolumn{1}{l}{compu} & Tenencia computadora & \multicolumn{1}{c}{1,57} & \multicolumn{1}{c}{0,49} & \multicolumn{1}{c}{1} & \multicolumn{1}{c}{2} & \multicolumn{1}{c}{13582} \\
    \multicolumn{1}{l}{casa} & Tenencia casa propia & \multicolumn{1}{c}{1,30} & \multicolumn{1}{c}{0,46} & \multicolumn{1}{c}{1} & \multicolumn{1}{c}{2} & \multicolumn{1}{c}{13551} \\
    \multicolumn{1}{l}{clase} & Autopercepción de clase & \multicolumn{1}{c}{3,73} & \multicolumn{1}{c}{0,96} & \multicolumn{1}{c}{1} & \multicolumn{1}{c}{5} &  \\
    \multicolumn{1}{l}{edad} & Edad  & \multicolumn{1}{c}{40,84} & 
    \multicolumn{1}{c}{16,48} & \multicolumn{1}{c}{16} & \multicolumn{1}{c}{100} &  \\
    \multicolumn{1}{l}{remoto} & Trabajo remoto  & \multicolumn{1}{c}{1,65} & 
    \multicolumn{1}{c}{0,47} & \multicolumn{1}{c}{1} & \multicolumn{1}{c}{2} &  \multicolumn{1}{c}{12927}\\
    \multicolumn{1}{l}{avg\_d\_kbps} & Calidad internet  & \multicolumn{1}{c}{30904} & 
    \multicolumn{1}{c}{32568} & \multicolumn{1}{c}{46} & \multicolumn{1}{c}{259191} &  \multicolumn{1}{c}{}\\
    \midrule
         N &      &  &       &       &       & 19749  \\
    
    \floatfoot{Nota: Se presentan las medias, desvíos estándar y valores máximos y mínimos para las variables utilizadas. Para las variables dicotómicas se presenta también la cantidad total de valores positivos. Se observa así, por ejemplo, que entre las 19749 personas encuestadas, 15680 hacen uso de algún tipo de red social pero solo 9618 poseen \textit{smartphone}.}
    \end{tabular}%
  \label{tab:Cuadro1}%
\end{table}%